\begin{table}[h]
    \centering
    \footnotesize
    % \setlength{\tabcolsep}{2.5pt}
    \caption{\textbf{受控本地部署下 LoCoMo 的详细延迟。}
    批大小为 1，所有数值的单位均为毫秒。
    \textbf{离线}表示平均离线记忆构建时间，
    \textbf{总计}表示平均总推理延迟，
    \textbf{过滤}表示平均检索/过滤模块延迟，
    \textbf{GPU} 表示平均 GPU 时间。
    N/A 表示不适用或未测量。}
    \label{tab:latency_locomo_detail}
    \begin{tabular}{@{}lrrrr@{}}
        \toprule
        \textbf{方法} & \textbf{离线} & \textbf{总计} & \textbf{过滤} & \textbf{GPU} \\
        \midrule
        MemoryOS & 40657 & 2842 & 1488 & N/A \\
        A-MEM    & 26842 & 1449 & 49   & N/A \\
        LightMem & 9740  & 1662 & 62   & N/A \\
        \midrule
        本文-IMP ($\lambda=0$)   & N/A & 3881 & 1176 & 46 \\
        本文-IMP ($\lambda=0.3$) & N/A & 2432 & 556  & 46 \\
        本文-IMP ($\lambda=0.9$) & N/A & 1167 & 46   & 46 \\
        \midrule
        本文-REA ($\lambda=0$)   & N/A & 3678 & 1274 & 47 \\
        本文-REA ($\lambda=0.3$) & N/A & 3429 & 1216 & 45 \\
        本文-REA ($\lambda=0.9$) & N/A & 3318 & 1209 & 48 \\
        \midrule
        本文-CAP ($\lambda=0$)   & N/A & 3461 & 1228 & 44 \\
        本文-CAP ($\lambda=0.3$) & N/A & 3058 & 1079 & 46 \\
        本文-CAP ($\lambda=0.9$) & N/A & 2853 & 976  & 52 \\
        \bottomrule
    \end{tabular}
\end{table}


\begin{table}[h]
    \centering
    \footnotesize
    % \setlength{\tabcolsep}{3pt}
    \caption{\textbf{LoCoMo 上的批大小扩展。}
    我们报告不同批大小下的平均总推理延迟，所有数值单位均为毫秒。
    \textbf{BS} 表示批大小。
    BudgetMem 报告 $\lambda=0$ 的本文-REA。}
    \label{tab:latency_locomo_batch}
    \begin{tabular}{@{}lrrrr@{}}
        \toprule
        \textbf{BS} & \textbf{MemoryOS} & \textbf{A-MEM} & \textbf{LightMem} & \textbf{本文-REA} \\
        \midrule
        1  & 2842 & 1449 & 1662 & 3678 \\
        2  & 1633 & 845  & 1059 & 2295 \\
        4  & 974  & 518  & 775  & 1327 \\
        8  & 604  & 296  & 438  & 858  \\
        16 & 387  & 154  & 269  & 547  \\
        32 & 203  & 96   & 163  & 326  \\
        \bottomrule
    \end{tabular}
\end{table}
